% =====================================================================
%  Chapter 4 — The Second Moment Method
%  File: chapters/04_second_moment.tex
% =====================================================================
\providecommand{\E}{\mathbb{E}}
\providecommand{\Prob}{\mathbb{P}}
\providecommand{\Var}{\operatorname{Var}}
\providecommand{\Cov}{\operatorname{Cov}}
\providecommand{\whp}{\textit{whp}}

\chapter{The Second Moment Method}
\label{chap:second-moment}

The first moment gives the mean of a random variable but nothing about its
spread. The \emph{second moment method} adds that control: if the variance of
$X$ is small next to $\E[X]^2$, then $X$ rarely strays far from its mean, and in
particular is usually positive. Throughout, \whp\ (``with high probability'')
means with probability tending to $1$ as $n\to\infty$.

\begin{proposition}{The second-moment criterion}{criterion}
For a nonnegative random variable $X$ with $\E[X]>0$,
\[
  \Prob(X=0)\le\frac{\Var(X)}{\E[X]^2}.
\]
Consequently, if $\E[X]>0$ and $\Var(X)=o\bigl(\E[X]^2\bigr)$ along a sequence,
then $X>0$ \whp.
\end{proposition}

\begin{proof}
If $X=0$ then $|X-\E X|=\E X$, so $\{X=0\}\subseteq\{|X-\E X|\ge\E X\}$, and
Chebyshev's inequality (Appendix~\ref{chap:toolkit}, Theorem~\ref{thm:chebyshev})
bounds the latter by $\Var(X)/\E[X]^2$.
\end{proof}

\begin{insight}
A second-moment proof is really a claim that the summands are nearly
uncorrelated. The mean $\E[X]=\sum_i\E[X_i]$ ignores dependence; the variance
$\Var(X)=\sum_i\Var(X_i)+\sum_{i\ne j}\Cov(X_i,X_j)$ measures it. So the work is
always to check that the covariance sum is small next to $\E[X]^2$. When it is,
we can pass from ``$\E[X]>0$'' to ``$X>0$ \whp'', which is what locating a
threshold requires.
\end{insight}

% =====================================================================
\section{Does a random graph contain a triangle?}
% =====================================================================

The \emph{Erd\H{o}s--R\'enyi random graph} $G(n,p)$ has $n$ vertices, each of the
$\binom n2$ edges present independently with probability $p$ (Appendix~\ref{chap:toolkit}).

\begin{question}{Triangle threshold}{triangle-q}
For which $p=p(n)$ does $G(n,p)$ contain a triangle \whp?
\end{question}

Let $X$ count the triangles. There are $\binom n3$ candidate triples, each a
triangle with probability $p^3$, so $\E[X]=\binom n3 p^3\asymp (np)^3$. This
suggests the transition happens near $p\asymp 1/n$, and the two moments confirm
it from opposite sides.

\begin{theorem}{Triangle threshold}{triangle-threshold}
If $np\to0$ then $G(n,p)$ is triangle-free \whp; if $np\to\infty$ then $G(n,p)$
contains a triangle \whp. Thus $p=1/n$ is a threshold.
\end{theorem}

\begin{proof}
If $np\to0$ then $\E[X]\asymp(np)^3\to0$, and Markov gives
$\Prob(X\ge1)\le\E[X]\to0$: no triangle \whp.

If $np\to\infty$ we use the second moment. Write $X=\sum_S\ind_S$ over triples
$S$. Two triangle-indicators are dependent only when their triples share an
edge, i.e.\ share two vertices. Grouping the covariance sum by overlap, the
only nonzero contributions come from pairs sharing exactly one edge (four
vertices, five edges), of which there are $\Theta(n^4)$, each contributing
$\Theta(p^5)$. Hence
\[
  \Var(X)\;\lesssim\;\E[X] + n^4 p^5\;\asymp\;n^3p^3+n^4p^5,
\]
while $\E[X]^2\asymp n^6p^6$. Therefore
\[
  \frac{\Var(X)}{\E[X]^2}\;\lesssim\;\frac{1}{(np)^3}+\frac{1}{n^2p}
  \;=\;\frac1{(np)^3}+\frac1{n\cdot(np)}\;\to\;0,
\]
so $X>0$ \whp\ by Proposition~\ref{prop:criterion}.
\end{proof}

\begin{insight}
The two moments split the work at the threshold. Below it $\E[X]\to0$, so Markov
alone kills the triangles. Above it $\E[X]\to\infty$ only says triangles are
expected on average; the variance bound is what rules out their all appearing
together on a rare event, and it is small here precisely because two random
triangles seldom share an edge. Right at $np\to c$ neither moment decides the
question, and the triangle count converges to a Poisson limit --- the threshold
is coarse, turning on over a constant factor rather than sharply.
\end{insight}

% =====================================================================
\section{Thresholds for fixed subgraphs}
% =====================================================================

The triangle argument generalizes to any fixed graph $H$, but with a twist: the
threshold is governed not by $H$ itself but by its densest subgraph.

\begin{definition}{Edge density and $m(H)$}{density}
The \emph{edge--vertex ratio} of a graph $H$ is $\rho(H)=e_H/v_H$ (half the
average degree), and
\[
  m(H)=\max_{H'\subseteq H}\rho(H').
\]
\end{definition}

The proof uses a packaged form of the second-moment estimate.

\begin{lemma}{Second-moment criterion via overlaps}{delta}
Let $X=\sum_i\ind_{A_i}$ with $i\sim j$ meaning $A_i,A_j$ are dependent, and set
$\Delta^\ast=\max_i\sum_{j\sim i}\Prob(A_j\mid A_i)$. If $\E[X]\to\infty$ and
$\Delta^\ast=o(\E[X])$, then $X>0$ \whp.
\end{lemma}

\begin{proof}
Since $\Var(X)\le\E[X]+\sum_i\Prob(A_i)\sum_{j\sim i}\Prob(A_j\mid A_i)
\le\E[X]\,(1+\Delta^\ast)$, Proposition~\ref{prop:criterion} gives
$\Prob(X=0)\le(1+\Delta^\ast)/\E[X]\to0$.
\end{proof}

\begin{theorem}{Subgraph threshold; Bollob\'as 1981}{subgraph-threshold}
For any fixed graph $H$, the function $p=n^{-1/m(H)}$ is a threshold for $G(n,p)$
to contain a copy of $H$.
\end{theorem}

\begin{proof}[Proof sketch]
Let $H'\subseteq H$ attain $\rho(H')=m(H)$. Below the threshold,
$\E[X_{H'}]\asymp n^{v_{H'}}p^{e_{H'}}=o(1)$, so $G(n,p)$ has no copy of $H'$,
hence none of $H$. Above the threshold, $\E[X_H]\to\infty$, and one checks
$\Delta^\ast=o(\E[X_H])$: the dominant overlaps between two copies of $H$
occur along a common subgraph $J$, and the condition $p\gg n^{-1/m(H)}\ge
n^{-1/\rho(J)}$ makes every such overlap term negligible against $\E[X_H]$.
Lemma~\ref{lem:delta} finishes the proof.
\end{proof}

\begin{insight}
A random graph gets $H$ as soon as it gets $H$'s densest subgraph; a sparse
piece hanging off a dense core shows up earlier and plays no role. That is why
the threshold is $n^{-1/m(H)}$ rather than $n^{-1/\rho(H)}$: for a pendant edge
on $K_4$ one has $\rho=7/5$ but $m=\rho(K_4)=3/2$, and it is the $K_4$ the graph
waits for.
\end{insight}

% =====================================================================
\section{The existence of thresholds}
% =====================================================================

Both examples above have a threshold. Is this a coincidence? For \emph{monotone}
properties it never is.

\begin{definition}{Threshold}{threshold}
A property $\mathcal P$ of $G(n,p)$ is \emph{increasing} if adding edges preserves
it. A sequence $r_n$ is a \emph{threshold} for $\mathcal P$ if $\Prob(G(n,p)\in
\mathcal P)\to0$ when $p/r_n\to0$ and $\to1$ when $p/r_n\to\infty$.
\end{definition}

\begin{theorem}{Bollob\'as--Thomason 1987}{threshold-exists}
Every nontrivial increasing property has a threshold.
\end{theorem}

\begin{proof}[Proof idea]
Monotonicity makes $p\mapsto\Prob(G(n,p)\in\mathcal P)$ strictly increasing, so a
\emph{critical} probability $p_c$ with $\Prob(G(n,p_c)\in\mathcal P)=1/2$ is
well defined. A \emph{multiple-exposure} argument --- running $m$ independent
copies of $G(n,p_c)$ and taking their union, which stochastically dominates
$G(n,mp_c)$ --- shows the probability jumps from near $0$ to near $1$ as $p$ moves
from $p_c/m$ to $mp_c$. Hence $p_c$ itself is a threshold.
\end{proof}

\begin{insight}
The theorem gives a threshold but not its width. Some properties ---
connectivity, or containing a fixed subgraph --- are sharp, with a transition
window narrower than the threshold; others, like containing a triangle at
$p=c/n$, are coarse, spread over a constant factor. Telling the two apart is
what the Friedgut--Bourgain theory of sharp thresholds is about.
\end{insight}

% =====================================================================
\section{The clique number of \texorpdfstring{$G(n,1/2)$}{G(n,1/2)}}
% =====================================================================

Let $\omega(G)$ be the clique number. For $G(n,1/2)$ the second moment pins it
down to two values. Write $f(n,k)=\binom nk 2^{-\binom k2}$, the expected number
of $k$-cliques.

\begin{theorem}{Second-moment bound for the clique number}{clique-sm}
If $f(n,k)\to0$ then $\omega(G(n,1/2))<k$ \whp; if $f(n,k)\to\infty$ then
$\omega(G(n,1/2))\ge k$ \whp.
\end{theorem}

\begin{proof}[Proof sketch]
The first claim is Markov: $\Prob(\omega\ge k)\le\E[X]=f(n,k)\to0$, where $X$
counts $k$-cliques. For the second, apply Lemma~\ref{lem:delta}: two $k$-sets
$S,T$ are dependent only if $|S\cap T|\ge2$, and summing
$\Prob(A_T\mid A_S)=2^{\binom{|S\cap T|}2-\binom k2}$ over such $T$ gives
$\Delta^\ast\ll\E[X]$ when $f(n,k)\to\infty$.
\end{proof}

\begin{theorem}{Two-point concentration; Bollob\'as--Erd\H{o}s, Matula 1976}{clique-2pt}
There is $k=k(n)\sim 2\log_2 n$ with
$\omega(G(n,1/2))\in\{k,k+1\}$ \whp.
\end{theorem}

\begin{proof}
The ratio of consecutive expectations is
\[
  \frac{f(n,k+1)}{f(n,k)}=\frac{n-k}{k+1}\,2^{-k}=n^{-1+o(1)}
  \qquad(k\sim2\log_2 n),
\]
so $f(n,k)$ plunges through every fixed level in a single step of $k$. Choose
$k_0\sim2\log_2 n$ with $f(n,k_0)\ge n^{-1/2}>f(n,k_0+1)$. Then
$f(n,k_0-1)\to\infty$ and $f(n,k_0+1)\to0$, so by Theorem~\ref{thm:clique-sm}
the clique number is confined to $\{k_0-1,k_0\}$ \whp.
\end{proof}

\begin{insight}
Two-point concentration is a strong conclusion: an optimization over
exponentially many subsets lands on one of two consecutive integers near
$2\log_2 n$. It works because $f(n,k)$ falls by a factor of about $n$ each time
$k$ grows by one, so there is no room in between. The same sharpness makes
$G(n,1/2)$ a convenient source of extremal objects --- its clique and independence
numbers are both $\sim2\log_2 n$ almost surely --- and Chapter~5 uses the
independence number to break the Haj\'os conjecture.
\end{insight}

% =====================================================================
\section{Hardy--Ramanujan and Erd\H{o}s--Kac}
% =====================================================================

The second moment reaches beyond graphs. Let $\nu(m)$ be the number of distinct
prime divisors of $m$.

\begin{theorem}{Hardy--Ramanujan 1917}{hardy-ramanujan}
For all but $o(n)$ of the integers $m\le n$, $\;\nu(m)=(1+o(1))\log\log n$.
\end{theorem}

\begin{proof}[Tur\'an's second-moment proof]
Pick $m$ uniformly from $[n]$ and write $\nu(m)\approx\sum_{p\le n}X_p$ where
$X_p=\ind_{p\mid m}$. Then $\E[X_p]=\lfloor n/p\rfloor/n\approx1/p$, so by
Mertens' theorem $\E[\nu]\approx\sum_{p\le n}1/p\approx\log\log n$. The pairs
$X_p,X_q$ are nearly uncorrelated ($\E[X_pX_q]\approx1/(pq)$ up to negligible
error), so $\Var(\nu)\approx\sum_p\tfrac1p(1-\tfrac1p)\approx\log\log n$ as well.
Chebyshev then gives $\nu(m)=(1+o(1))\log\log n$ outside an $o(n)$-sized
exceptional set.
\end{proof}

\begin{insight}
The heuristic is that ``divisible by $p$'' behaves like an independent coin of
bias $1/p$ across the primes. The second moment makes this precise enough to get
concentration; the method of moments pushes it to a central limit theorem
(Erd\H{o}s--Kac): $\bigl(\nu(m)-\log\log n\bigr)/\sqrt{\log\log n}$ tends to a
standard normal as $m$ ranges over $[n]$. That result is where probabilistic
number theory began.
\end{insight}

% =====================================================================
\section{Sets with distinct subset sums}
% =====================================================================

\begin{question}{Distinct sums}{distinct-q}
How small can the largest element of a set of $n$ positive integers be, if all
$2^n$ subset sums are distinct? (Powers of two give largest element $2^{n-1}$;
can one do better?)
\end{question}

\begin{theorem}{Erd\H{o}s's lower bound}{distinct-sums}
If $a_1<\cdots<a_n$ are positive integers with distinct subset sums, then
$a_n\ge c\,2^n/\sqrt n$ for an absolute constant $c>0$.
\end{theorem}

\begin{proof}
Let $\varepsilon_1,\dots,\varepsilon_n$ be independent uniform signs and put
$X=\sum_i\varepsilon_i a_i$. Then $\E[X]=0$ and
$\Var(X)=\sum_i a_i^2\le n\,a_n^2$. By Chebyshev, at least $\tfrac34$ of the
$2^n$ sign patterns yield $|X|\le 2\sqrt{\Var(X)}\le 2\sqrt n\,a_n$, i.e.\ lie
in an interval of length $4\sqrt n\,a_n$. But distinct subset sums force the
$2^n$ values of $X$ to be distinct integers, and an interval of length $L$
holds at most $L+1$ of them. Hence $\tfrac34\cdot2^n\le 4\sqrt n\,a_n+1$, giving
$a_n\ge c\,2^n/\sqrt n$.
\end{proof}

\begin{insight}
The bound comes from the spread of a random signed sum, not from any
construction. Chebyshev forces the $2^n$ subset sums into an interval of width
$O(\sqrt n\,a_n)$; distinctness lets only $O(\sqrt n\,a_n)$ integers sit there;
comparing the two makes $a_n$ exponential. Erd\H{o}s offered money for closing
the $\sqrt n$ gap to the conjectured $a_n\ge c\,2^n$; Dubroff, Fox, and Xu
(2021) reached $2^n/\sqrt n$ with the best known constant, using Harper's
vertex-isoperimetric inequality on the cube.
\end{insight}

% =====================================================================
\section{The Weierstrass approximation theorem}
% =====================================================================

The method is not confined to combinatorics; here it proves a classical
analysis theorem.

\begin{theorem}{Weierstrass 1885}{weierstrass}
Every continuous $f\colon[0,1]\to\mathbb R$ is a uniform limit of polynomials.
\end{theorem}

\begin{proof}
Define the Bernstein polynomials
$B_n(x)=\sum_{k=0}^n f\!\bigl(\tfrac kn\bigr)\binom nk x^k(1-x)^{n-k}$. If
$S\sim\mathrm{Binomial}(n,x)$ then $B_n(x)=\E\bigl[f(S/n)\bigr]$, and
$S/n$ concentrates at $x$ because $\Var(S/n)=x(1-x)/n\le 1/(4n)$. Fix
$\varepsilon>0$; uniform continuity gives $\delta$ with
$|f(u)-f(v)|<\varepsilon$ when $|u-v|<\delta$, and $\|f\|_\infty=M<\infty$.
Splitting on whether $|S/n-x|<\delta$,
\[
  |B_n(x)-f(x)|\le\E\bigl|f(S/n)-f(x)\bigr|
  \le\varepsilon+2M\,\Prob\bigl(|S/n-x|\ge\delta\bigr)
  \le\varepsilon+\frac{2M}{4n\delta^2},
\]
using Chebyshev. The bound is uniform in $x$ and $<2\varepsilon$ for $n$ large,
so $B_n\to f$ uniformly.
\end{proof}

\begin{insight}
A Bernstein polynomial is an expectation: flip $n$ coins of bias $x$, see where
the average lands, and read off $f$ there. Since $S/n$ concentrates at $x$, its
average value of $f$ sits near $f(x)$. The point worth keeping is that the same
variance estimate used for random graphs also settles a theorem in
approximation theory.
\end{insight}

% =====================================================================
\section*{Exercises}
\addcontentsline{toc}{section}{Exercises}
% =====================================================================

\begin{enumerate}[itemsep=4pt]
  \item Let $X\ge0$ with $\Prob(X=0)<1$. Prove
        $\Prob(X=0)\le\Var(X)/\E[X^2]$.
  \item For $X$ with mean $\mu$ and variance $\sigma^2$, prove the one-sided
        bound $\Prob(X\ge\mu+\lambda)\le\sigma^2/(\sigma^2+\lambda^2)$ for all
        $\lambda>0$.
  \item (Threshold for $k$-APs.) For $[n]_p$ the random subset keeping each
        element with probability $p$, determine the threshold for containing a
        $k$-term arithmetic progression, for each fixed $k\ge3$.
  \item Determine the threshold function for $G(n,p)$ to contain a cycle.
  \item Show that for each fixed $k$ there is a sequence $p_n$ with
        $\Prob\bigl(G(n,p_n)\text{ has a component of exactly }k\text{ vertices}\bigr)\to1$.
  \item (Poisson limit.) For $X$ the number of triangles in $G(n,c/n)$, compute
        $\lim_n\E\binom Xk$ for each $k$, match it to
        $\mathrm{Binomial}(n,\lambda/n)$, and deduce the limiting distribution of
        $X$ and of $\Prob(X>0)$.
  \item (CLT for triangle counts.) With $X$ the triangle count and $Y$ the edge
        count of $G(n,1/2)$, find $a_n$ with $\Var(X-a_nY)=o(\Var X)$, and deduce
        that $X$ is asymptotically normal.
\end{enumerate}
